\section{Structured Batching for Interleaved Codeword}
\label{sec:interleave}

In this section,
we discuss code switching from an interleaved code $\Code^t: \Field^{k\times t} \mapsto \Field^{\mu \times t}$ to a second code $\Code':\Field^{k} \mapsto \Field^{\mu'}$.
This requires batching the interleaved code $\bm{C} \in \Field^{\mu \times t}$ to $\bm{u} = \bm{C}\bm{r}$ using some random $\bm{r} \in \Field^t$ before code switching.
The proximity between $\bm{u}$ and $\bm{C}\bm{r}$ are ensured by checking random $\ell$ rows of $\bm{C}$, denoted by $\bm{C}_\mathcal{I}$ with the selected index set $\mathcal{I}$.
This idea was first proposed by Blaze~\cite{eurocrypt/BrehmCFRRZ25}.
Unlike Blaze, we use a tensor-span vector $\bm{r} \in \TensorSpace(t)$ for batching.
The advantage is that the proximity check is reduced to checking an MLE evaluation of $\bm{C}_\mathcal{I}$,
allowing it to be verified via another MLE oracle.




\subsection{MLE Oracle for Interleaved Code}
\label{sec:interleaveOracle}

Form Lemma~\ref{lemma:odd},
to obtain an MLE oracle for an interleaved code $\Code^t$,
we just need to design and MLE-ODD-IOPP for the relation in Eq.~\ref{eq:relationMLEoddiopp}.

Such an IOPP can achieved by structured batching the codeword $\bm{C}\in \Field^{\mu \times t}$
into a single column vector $\bm{u}:=\bm{C}\tilde{\bm{r}} \in \Field^\mu$,
and check the proximity of $\bm{u}$ to the codeword of $\bm{M}\tilde{\bm{r}}$.

Following this idea, we present the protocol of such MLE-ODD-IOPP in Algorithm~\ref{alg:interleave}
with the following security guarantees.

We ensure the proximity such that $\Code^t$ is list decodable:
\[
    \delta^{\odot t}(\bm{C}, \Enc[\Code^t](\bm{M})) \le \epsilon,
\]

% Here, $\delta_\mathsf{row}$ denotes the row-wise relative distance, meaning that $\bm{C}$ and $\Enc[\Code^t](\bm{M})$ differ in at most $\epsilon \cdot \mu$ rows.
% The  distance is at least the relative Hamming distance:
% \[
%     \delta^{\odot t}(\bm{C}, \Enc[\Code^t](\bm{M}))
%     \ge \delta_H(\bm{C}, \Enc[\Code^t](\bm{M})),
% \]
% since a single disagreeing row is counted as one full unit in the row-wise distance.

\begin{theorem}[Interleaved MLE-ODD-IOPP]
\label{theorem:interleavedIOPP}
Algorithm~\ref{alg:interleave} functions as an IOPP for the following relation:
\begin{equation*}
% \label{eq:relationMLEoddiopp}
\Relation_{\mathsf{MLE\mbox{-}ODD\mbox{-}IOPP}}[\epsilon]
=
\RelationExp{
\begin{aligned}
    &\Code,\, \bm{C},\, \epsilon \\
    &\hat{v}, \hat{\bm{l}}, \hat{\bm{r}}, v, \bm{l}, \bm{r}
\end{aligned}
}{\bm{M}}{
\begin{aligned}
    &\hat{v} = \hat{\bm{l}}^{\top}\bm{M}\hat{\bm{r}}
    \;\wedge\;
    v = \bm{l}^{\top}\bm{M}\bm{r} \\
    &\wedge\;
    \delta^{\otimes t}(\bm{C}, \Enc[\Code^t](\bm{M})) \le \epsilon
\end{aligned}
}
\end{equation*}
with computational soundness error and knowledge soundness error bounded by:
\begin{equation}
\label{eq:interleaveError}
      (1-\epsilon)^\ell + \frac{(\mu+2)\log_2 t + 2 \log_2 \ell +3}{|\Field|-2}.
\end{equation}
\end{theorem}

Proof of Theorem~\ref{theorem:interleavedIOPP} is provided in Appendix~\ref{appx:proofInterleavedIOPP}.

\input{alg/alg-interleave}

Combining Theorem~\ref{theorem:interleavedIOPP} with Lemma~\ref{lemma:odd} and Theorem~\ref{theorem:iop},
we obtain an MLE oracle for interleaved codes by modifying the protocol of Algorithm~\ref{alg:interleave} as follows.
\begin{itemize}
    \item First, in Line~\ref{line:interleaveQueryC}, instead of directly retrieving the entries of $\bm{C}_{\mathcal{I}}$,
the verifier issues point-oracle queries to $\Oracle{\mathsf{P}}(\bm{C})$.
    \item Second, in Line~\ref{line:interleaveSendv}, rather than sending the vector $\bm{v}$ explicitly, the prover supplies an MLE oracle for a suitable auxiliary code, $\Oracle_{\mathsf{MLE}}[\Code']( \bm{v} )$.
    \item Accordingly, the verification check in Line~\ref{line:interleaveCheckV} is replaced by matrix-function queries to the MLE oracle $\Oracle_{\mathsf{MLE}[\Code']}( \bm{v} )$.
\end{itemize}

As we will see in the next section, choosing different pairs of codes
$\Code$ and $\Code'$ yields a wide range of MLE oracles for product codes.

\medskip
\noindent
{\bf Query complexity.}
In the soundness bound in Eq.~\eqref{eq:interleaveError},
the last term depends on the codeword dimension, which is proportional to the witness size.
By choosing a field of size at least $2^\lambda \cdot \mu \cdot \log_2 t$,
this term can be bounded by $2^{-\lambda}$.

To bound the first term, it suffices to require:
\[
    (1-\delta_\mathsf{LD})^{\ell} \le 2^{-\lambda}
    \quad\Longleftrightarrow\quad
    \ell \ge \frac{\lambda}{-\log_2(1-\delta_\mathsf{LD})}.
\]
For the rate-$1/4$ RAA code, $\delta = 0.19$~\cite{eurocrypt/BrehmCFRRZ25} and $\delta_\mathsf{LD} \ge 0.19/2$,
which yields:
\[
    \ell
    = \left\lceil \frac{1}{-\log_2(1-0.19/2)} \right\rceil \cdot \lambda
    = 7\,\lambda.
\]
Therefore, setting $\ell = 7\lambda$ suffices to achieve $2^{-\lambda}$ security.

Since each row of $\bm{C}_{\mathcal{I}}$ consists of $t$ entries,
the verifier must first compute the Merkle root of these $t$ entries using $t$ hash evaluations,
and then verify the resulting row Merkle root using $\log_2 \mu$ additional hash evaluations along the Merkle path.
Thus, by choosing $t \le \log_2(\mu t^2)$, \eg, $t=32$ when $kt = 2^{30}$ and using a rate-$1/4$ RAA code,
the verifier performs $(\log_2 \mu + t)\cdot 7\lambda \le 2\log_2 (\mu t)\cdot 7\lambda$ hash evaluations, less than $14\lambda$ full Merkle authentication paths in $\Oracle_\mathsf{P}(\bm{C})$.

\medskip
\noindent
\textbf{Comparison with Blaze~\cite{eurocrypt/BrehmCFRRZ25}.}
Our construction of interleaved-code MLE oracles is conceptually related to
Blaze~\cite{eurocrypt/BrehmCFRRZ25}, but differs in several key aspects:
\begin{itemize}
    \item We use a \emph{structured} batching vector
    $\tilde{\bm{r}} \in \TensorSpace(t)$,
    whereas Blaze samples an unstructured vector
    $\tilde{\bm{r}} \in \Field^t$.
    \item Blaze requires the proximity parameter to satisfy
    $\epsilon \le \delta/3$~\cite[Lemma~6.1]{eurocrypt/BrehmCFRRZ25},
    while our construction only requires
    $\epsilon \le \delta_{\mathsf{LD}}$ due to the use of OOD sampling.
    % \item As a consequence, our approach supports product codes as the base code $\Code$,
    % even when their unique decoding radius is small, provided their list-decoding
    % radius is sufficiently large.
\end{itemize}

Due to these differences, we refer to our interleaved MLE oracle as
\blazeplus.

\iffalse
\medskip
\noindent
{\bf Interleaved Code as the Base Code.}
In the next section, we will instantiate $\Code$ in Algorithm~\ref{alg:interleave}
as the interleaved code $\Code_{\mathsf{base}}^{\gamma}$ for some
$\gamma \in \mathbb{N}$.
In this case, the generator matrix of $\Code$, denoted by $\bm{G}$,
can be written as:
\[
    \bm{G} = \bm{G}_{\mathsf{base}} \kronecker \bm{I}_{\gamma},
\]
where $\bm{G}_{\mathsf{base}}$ is the generator matrix of
$\Code_{\mathsf{base}}$ and $\bm{I}_{\gamma}$ is the
$\gamma \times \gamma$ identity matrix.
By Corollary~\ref{corollary:mleGeneratorInterleave},
this generator matrix admits efficient code switching.
\fi

